\newgeometry{margin=.5cm}
\pagestyle{empty}
\raggedcolumns
\footnotesize
\newcommand{\cheatSheetColor}{RoyalBlue}

\newtcolorbox{cheatSheetBox}[1][]{
	enhanced jigsaw,
	colback=\cheatSheetColor!10,
	colframe=\cheatSheetColor,
	colbacktitle=\cheatSheetColor!30,
	fonttitle=\bfseries,
	parbox=false,
	attach boxed title to top text left={yshift=-3mm,yshifttext=-3mm},
	coltitle=black,
	#1}

\begin{landscape}
\footnotesize
\begin{multicols*}{4}

\textcolor{\cheatSheetColor}{
\begin{center}
\section{Statistik Cheatsheet}
\end{center}\hrule}

\textcolor{\cheatSheetColor}{\subsection*{Skalenniveaus}}

\begin{cheatSheetBox}[title=Nominalskala]
Kategorien ohne Ordnung.

Beispiele:
\begin{itemize}
\item Geschlecht
\item Nationalität
\item Lieblingsmusik
\end{itemize}

Sinnvolle Kennzahlen:
\begin{itemize}
\item Häufigkeit
\item Modus
\end{itemize}
\end{cheatSheetBox}

\begin{cheatSheetBox}[title=Ordinalskala]
Kategorien mit Rangordnung.

Beispiele:
\begin{itemize}
\item Schulnoten
\item Ranglisten
\item Zufriedenheitsskala
\end{itemize}

Sinnvolle Kennzahlen:
\begin{itemize}
\item Median
\item Quartile
\end{itemize}
\end{cheatSheetBox}

\begin{cheatSheetBox}[title=Intervallskala]
Gleiche Abstände, aber kein echter Nullpunkt.

Beispiel:
\begin{itemize}
\item Temperatur in °C
\end{itemize}

Erlaubt:
\begin{itemize}
\item Mittelwert
\item Standardabweichung
\end{itemize}
\end{cheatSheetBox}

\begin{cheatSheetBox}[title=Verhältnisskala]
Numerische Werte mit echtem Nullpunkt.

Beispiele:
\begin{itemize}
\item Gewicht
\item Länge
\item Einkommen
\end{itemize}

Verhältnisse sind sinnvoll.
\end{cheatSheetBox}

\textcolor{\cheatSheetColor}{\subsection*{Lagekennzahlen}}

\begin{cheatSheetBox}[title=Minimum und Maximum]
Kleinster und grösster Wert im Datensatz.

Spannweite:
\[
\text{Spannweite} = \text{Max} - \text{Min}
\]
\end{cheatSheetBox}

\begin{cheatSheetBox}[title=Mittelwert]
Durchschnitt aller Werte.

\[
\bar{x} = \frac{1}{n}\sum_{i=1}^{n} x_i
\]

Empfindlich gegenüber Ausreissern.
\end{cheatSheetBox}

\begin{cheatSheetBox}[title=Median]
Mittlerer Wert nach Sortierung.

Robust gegenüber Ausreissern.

Beispiel:

1, 2, 3, 4, 100

Median = 3
\end{cheatSheetBox}

\textcolor{\cheatSheetColor}{\subsection*{Streuungsmasse}}

\begin{cheatSheetBox}[title=Quartile]
Datensatz wird in vier Teile geteilt.

\begin{itemize}
\item Q1 = 25\%
\item Median = 50\%
\item Q3 = 75\%
\end{itemize}
\end{cheatSheetBox}

\begin{cheatSheetBox}[title=Quartilsabstand]
Interquartilsabstand:

\[
IQR = Q3 - Q1
\]

Robust gegenüber Ausreissern.
\end{cheatSheetBox}

\begin{cheatSheetBox}[title=Standardabweichung]

Misst Streuung um den Mittelwert.

\[
s = \sqrt{\frac{1}{n-1}\sum (x_i-\bar{x})^2}
\]

Gross → Daten stark verteilt.
\end{cheatSheetBox}

\textcolor{\cheatSheetColor}{\subsection*{Darstellungen}}

\begin{cheatSheetBox}[title=Histogramm]
Zeigt Verteilung numerischer Daten.

Gut für:
\begin{itemize}
\item Form der Verteilung
\item Schiefe
\item Ausreisser
\end{itemize}
\end{cheatSheetBox}

\begin{cheatSheetBox}[title=Boxplot]

Zeigt:

\begin{itemize}
\item Median
\item Quartile
\item Ausreisser
\end{itemize}

Sehr kompakte Darstellung.
\end{cheatSheetBox}

\begin{cheatSheetBox}[title=Stem-and-Leaf Plot]

Gut für kleine Datensätze.

Beispiel:

\[
2 | 3\,5\,8
\]

→ 23, 25, 28
\end{cheatSheetBox}

\textcolor{\cheatSheetColor}{\subsection*{Typische Fehler}}

\begin{cheatSheetBox}[title=Statistik kritisch lesen]

Achten Sie auf:

\begin{itemize}
\item manipulierte Achsen
\item kleine Stichproben
\item Mittelwert statt Median
\item fehlende Vergleichsgruppen
\end{itemize}

Statistik kann leicht falsch interpretiert werden.
\end{cheatSheetBox}

\end{multicols*}
\end{landscape}